\begin{abcliste}
\abc The magnet turns twice as fast, so the flux through the coil changes twice as fast: the voltage is $\result{\text{twice as large, and it alternates twice as often}}$.

\abc When the wheel stands still, the magnet does not turn: the flux does not change, no voltage is induced, and the lamp is $\result{\text{off}}$.

\abc With the lamp on, a current flows in the coil. By Lenz's rule, it opposes the change that causes it, the turning of the magnet: pedalling is $\result{\text{harder}}$. The energy for the light comes from your legs.
\end{abcliste}
